\documentclass[10pt,a4paper]{amsart}
\usepackage{amsmath,amssymb,amsthm,graphicx}
\usepackage{verbatim}
\newenvironment{reference}{\comment}{\endcomment}
%\newenvironment{reference}{}{}
\newenvironment{slogan}{\comment}{\endcomment}
\newenvironment{history}{\comment}{\endcomment}

% For commutative diagrams we use Xy-pic
\usepackage[all]{xy}

% We use 2cell for 2-commutative diagrams.
\xyoption{2cell}
% Theorem environments.
%
\theoremstyle{plain}
\newtheorem{theorem}[subsection]{Theorem}
\newtheorem{proposition}[subsection]{Proposition}
\newtheorem{lemma}[subsection]{Lemma}

\theoremstyle{definition}
\newtheorem{definition}[subsection]{Definition}
\newtheorem{example}[subsection]{Example}
\newtheorem{exercise}[subsection]{Exercise}
\newtheorem{situation}[subsection]{Situation}

\theoremstyle{remark}
\newtheorem{remark}[subsection]{Remark}
\newtheorem{remarks}[subsection]{Remarks}

\numberwithin{equation}{subsection}

% Macros
%
\def\lim{\mathop{\mathrm{lim}}\nolimits}
\def\colim{\mathop{\mathrm{colim}}\nolimits}
\def\Spec{\mathop{\mathrm{Spec}}}
\def\Hom{\mathop{\mathrm{Hom}}\nolimits}
\def\Ext{\mathop{\mathrm{Ext}}\nolimits}
\def\SheafHom{\mathop{\mathcal{H}\!\mathit{om}}\nolimits}
\def\SheafExt{\mathop{\mathcal{E}\!\mathit{xt}}\nolimits}
\def\Sch{\mathit{Sch}}
\def\Mor{\mathop{\mathrm{Mor}}\nolimits}
\def\Ob{\mathop{\mathrm{Ob}}\nolimits}
\def\Sh{\mathop{\mathit{Sh}}\nolimits}
\def\NL{\mathop{N\!L}\nolimits}
\def\CH{\mathop{\mathrm{CH}}\nolimits}
\def\proetale{{pro\text{-}\acute{e}tale}}
\def\etale{{\acute{e}tale}}
\def\QCoh{\mathit{QCoh}}
\def\Ker{\mathop{\mathrm{Ker}}}
\def\Im{\mathop{\mathrm{Im}}}
\def\Coker{\mathop{\mathrm{Coker}}}
\def\Coim{\mathop{\mathrm{Coim}}}

% Boxtimes
%
\DeclareMathSymbol{\boxtimes}{\mathbin}{AMSa}{"02}

%
% Macros for moduli stacks/spaces
%
\def\QCohstack{\mathcal{QC}\!\mathit{oh}}
\def\Cohstack{\mathcal{C}\!\mathit{oh}}
\def\Spacesstack{\mathcal{S}\!\mathit{paces}}
\def\Quotfunctor{\mathrm{Quot}}
\def\Hilbfunctor{\mathrm{Hilb}}
\def\Curvesstack{\mathcal{C}\!\mathit{urves}}
\def\Polarizedstack{\mathcal{P}\!\mathit{olarized}}
\def\Complexesstack{\mathcal{C}\!\mathit{omplexes}}
% \Pic is the operator that assigns to X its picard group, usage \Pic(X)
% \Picardstack_{X/B} denotes the Picard stack of X over B
% \Picardfunctor_{X/B} denotes the Picard functor of X over B
\def\Pic{\mathop{\mathrm{Pic}}\nolimits}
\def\Picardstack{\mathcal{P}\!\mathit{ic}}
\def\Picardfunctor{\mathrm{Pic}}
\def\Deformationcategory{\mathcal{D}\!\mathit{ef}}
\usepackage{cleveref}
\crefname{originallemma}{original lemma}{original lemmas}
\Crefname{originallemma}{Original lemma}{Original lemmas}
\makeatletter
\newlabel{lemma-quotient-stack-functorial}{{046Q}{source}{}{sourcepointer-1}{}}
\newlabel{lemma-quotient-stack-functorial@cref}{{[originallemma][0][]046Q}{[1][0][]source}{}{}{}}
\makeatother
\crefname{originallemma}{original lemma}{original lemmas}
\Crefname{originallemma}{Original lemma}{Original lemmas}
\makeatletter
\newlabel{lemma-quotient-stack-restrict}{{046S}{source}{}{sourcepointer-2}{}}
\newlabel{lemma-quotient-stack-restrict@cref}{{[originallemma][0][]046S}{[1][0][]source}{}{}{}}
\makeatother
\crefname{originalreference}{original reference}{original references}
\Crefname{originalreference}{Original reference}{Original references}
\makeatletter
\newlabel{stacks-lemma-characterize-essentially-surjective-when-ff}{{046N}{source}{}{sourcepointer-3}{}}
\newlabel{stacks-lemma-characterize-essentially-surjective-when-ff@cref}{{[originalreference][0][]046N}{[1][0][]source}{}{}{}}
\makeatother
\crefname{originallemma}{original lemma}{original lemmas}
\Crefname{originallemma}{Original lemma}{Original lemmas}
\makeatletter
\newlabel{lemma-quotient-stack-objects}{{044X}{source}{}{sourcepointer-4}{}}
\newlabel{lemma-quotient-stack-objects@cref}{{[originallemma][0][]044X}{[1][0][]source}{}{}{}}
\makeatother
\crefname{originalreference}{original reference}{original references}
\Crefname{originalreference}{Original reference}{Original references}
\makeatletter
\newlabel{spaces-topologies-lemma-fppf-covering-surjective}{{0469}{source}{}{sourcepointer-5}{}}
\newlabel{spaces-topologies-lemma-fppf-covering-surjective@cref}{{[originalreference][0][]0469}{[1][0][]source}{}{}{}}
\makeatother
\title{Equivalence criterion for a quotient-stack morphism}
\author{The Stacks Project Authors}
\date{stacks-project:spaces-groupoids.tex:local-snapshot-20261012; GNU FDL 1.2 or later}
\begin{document}
\maketitle
\noindent\textit{Editorial note.} All delivered mathematical text is the author's own, copied verbatim from the pinned Stacks Project chapter file spaces-groupoids.tex; no statement or proof step was rewritten, reordered or removed.\par\medskip
\noindent\textit{Editorial note.} Span recomputation by label pairing (LANE\_BRIEF 10.2): in this chapter the author's proof environment follows the statement environment as a standalone block; the spec covers the complete environment pair (statement span, proof span) as recorded in the pool candidate.\par\medskip
\ifdefined\hypertarget\hypertarget{sourcepointer-1}{}\fi
\noindent\textit{Source reference.} Original lemma 046Q: sources/stacks-chow/tags.txt, source lines 5481--5481. Cross-chapter standard prerequisite cited by the author inside the counted span; it is used as a pointer and its own proof is not part of this item.\par\medskip
\ifdefined\hypertarget\hypertarget{sourcepointer-2}{}\fi
\noindent\textit{Source reference.} Original lemma 046S: sources/stacks-chow/tags.txt, source lines 5483--5483. Cross-chapter standard prerequisite cited by the author inside the counted span; it is used as a pointer and its own proof is not part of this item.\par\medskip
\ifdefined\hypertarget\hypertarget{sourcepointer-3}{}\fi
\noindent\textit{Source reference.} Original reference 046N: sources/stacks-chow/tags.txt, source lines 5478--5478. Cross-chapter standard prerequisite cited by the author inside the counted span; it is used as a pointer and its own proof is not part of this item.\par\medskip
\ifdefined\hypertarget\hypertarget{sourcepointer-4}{}\fi
\noindent\textit{Source reference.} Original lemma 044X: sources/stacks-chow/tags.txt, source lines 5416--5416. Cross-chapter standard prerequisite cited by the author inside the counted span; it is used as a pointer and its own proof is not part of this item.\par\medskip
\ifdefined\hypertarget\hypertarget{sourcepointer-5}{}\fi
\noindent\textit{Source reference.} Original reference 0469: sources/stacks-chow/tags.txt, source lines 5464--5464. Cross-chapter standard prerequisite cited by the author inside the counted span; it is used as a pointer and its own proof is not part of this item.\par\medskip
\begin{lemma}
\label{lemma-quotient-stack-restrict-equivalence}
Notation and assumption as in
Lemma \ref{lemma-quotient-stack-functorial}.
The morphism of quotient stacks
$$
[f] : [U/R] \longrightarrow [U'/R']
$$
is an equivalence if and only if
\begin{enumerate}
\item $(U, R, s, t, c)$ is the restriction of $(U', R', s', t', c')$
via $f : U \to U'$, and
\item the map
$$
\xymatrix{
U \times_{f, U', t'} R' \ar[r]_-{\text{pr}_1} \ar@/^3ex/[rr]^h
& R' \ar[r]_{s'} & U'
}
$$
is a surjection of sheaves.
\end{enumerate}
Part (2) holds for example if $\{h : U \times_{f, U', t'} R' \to U'\}$
is an fppf covering, or if $f : U \to U'$ is a surjection of sheaves, or if
$\{f : U \to U'\}$ is an fppf covering.
\end{lemma}

\begin{proof}
We already know that part (1) is equivalent to
fully faithfulness by
Lemma \ref{lemma-quotient-stack-restrict}.
Hence we may assume that (1) holds and that $[f]$ is fully faithful.
Our goal is to show, under these assumptions, that $[f]$ is an equivalence
if and only if (2) holds. We may use
Stacks, Lemma \ref{stacks-lemma-characterize-essentially-surjective-when-ff}
which characterizes equivalences.

\medskip\noindent
Assume (2). We will use
Stacks, Lemma \ref{stacks-lemma-characterize-essentially-surjective-when-ff}
to prove $[f]$ is an equivalence.
Suppose that $T$ is a scheme and $x' \in \Ob([U'/R']_T)$.
There exists a covering $\{g_i : T_i \to T\}$ such that $g_i^*x'$
is the image of some element $a'_i \in U'(T_i)$, see
Lemma \ref{lemma-quotient-stack-objects}.
Hence we may assume that $x'$ is the image of $a' \in U'(T)$.
By the assumption that $h$ is a surjection of sheaves, we
can find an fppf covering $\{\varphi_i : T_i \to T\}$ and morphisms
$b_i : T_i \to U \times_{g, U', t'} R'$ such that
$a' \circ \varphi_i = h \circ b_i$. Denote
$a_i = \text{pr}_0 \circ b_i : T_i \to U$. Then we see that
$a_i \in U(T_i)$ maps to $f \circ a_i \in U'(T_i)$ and
that $f \circ a_i \cong_{T_i} h \circ b_i = a' \circ \varphi_i$,
where $\cong_{T_i}$ denotes isomorphism in the fibre category
$[U'/R']_{T_i}$. Namely, the element of $R'(T_i)$ giving the isomorphism
is $\text{pr}_1 \circ b_i$. This means that the restriction
of $x$ to $T_i$ is in the essential image of the functor
$[U/R]_{T_i} \to [U'/R']_{T_i}$ as desired.

\medskip\noindent
Assume $[f]$ is an equivalence. Let $\xi' \in [U'/R']_{U'}$ denote the
object corresponding to the identity morphism of $U'$. Applying
Stacks, Lemma \ref{stacks-lemma-characterize-essentially-surjective-when-ff}
we see there exists an fppf covering $\mathcal{U}' = \{g'_i : U'_i \to U'\}$
such that $(g'_i)^*\xi' \cong [f](\xi_i)$ for some $\xi_i$ in $[U/R]_{U'_i}$.
After refining the covering $\mathcal{U}'$ (using
Lemma \ref{lemma-quotient-stack-objects})
we may assume $\xi_i$ comes from a morphism $a_i : U'_i \to U$.
The fact that $[f](\xi_i) \cong (g'_i)^*\xi'$ means that, after possibly
refining the covering $\mathcal{U}'$ once more, there exist morphisms
$r'_i : U'_i \to R'$ with $t' \circ r'_i = f \circ a_i$ and
$s' \circ r'_i = \text{id}_{U'} \circ g'_i$. Picture
$$
\xymatrix{
U \ar[d]^f & & U'_i \ar[ll]^{a_i} \ar[ld]_{r'_i} \ar[d]^{g'_i} \\
U' & R' \ar[l]_{t'} \ar[r]^{s'} & U'
}
$$
Thus $(a_i, r'_i) : U'_i \to U \times_{g, U', t'} R'$ are morphisms
such that $h \circ (a_i, r'_i) = g'_i$ and we conclude that
$\{h : U \times_{g, U', t'} R' \to U'\}$ can be refined by the
fppf covering $\mathcal{U}'$ which means that $h$ induces a surjection
of sheaves, see
Topologies on Spaces, Lemma
\ref{spaces-topologies-lemma-fppf-covering-surjective}.

\medskip\noindent
If $\{h\}$ is an fppf covering, then it induces a surjection of sheaves, see
Topologies on Spaces,
Lemma \ref{spaces-topologies-lemma-fppf-covering-surjective}.
If $U' \to U$ is surjective, then also $h$ is surjective as $s$ has a section
(namely the neutral element $e$ of the groupoid in algebraic spaces).
\end{proof}

\end{document}
